\chapter{Allergies and Hypersensitivities}
\index{Allergies and Hypersensitivities}

\section*{Definition and Overview}
Hypersensitivity is an umbrella term for harmful immune reactions in which the immune system damages the body's own tissues or produces symptoms in response to an otherwise harmless substance. Allergy is the best-known form of hypersensitivity, involving immunoglobulin E (IgE) and immediate symptoms, but hypersensitivity encompasses a wider range of mechanisms. The classical Gell and Coombs classification divides hypersensitivity reactions into four types, based on the immune components responsible and the time course of the reaction. Understanding the type is clinically important, because recognition, investigation, and treatment differ substantially between them. Hypersensitivity reactions can be provoked by foods, medicines, insect venom, contactants, and environmental agents, and they range from trivial skin irritation to life-threatening anaphylaxis or severe drug reactions.

\section*{Epidemiology}
Hypersensitivity conditions are extremely common. Allergic rhinitis affects around one in five people in many countries, and eczema, asthma, food allergy, and drug allergy affect substantial proportions of the population. Drug hypersensitivity reactions account for a significant minority of all adverse drug events and are a common cause of hospital morbidity. Autoimmune diseases, which share mechanisms with certain hypersensitivity reactions, become more frequent with age, whereas immediate-type allergies such as food allergy are more common in childhood. A family history of atopy is the strongest risk factor for the allergic forms, and the prevalence of allergic disease has risen steadily over recent decades, particularly in urbanised and industrialised societies.

\section*{Aetiology and Pathophysiology}
The immune system ordinarily protects against infection, but in hypersensitivity it responds inappropriately to a trigger. The four principal mechanisms are:

\begin{itemize}
    \item \textbf{Type I (immediate, IgE-mediated):} allergen-specific IgE bound to mast cells and basophils is cross-linked on re-exposure, releasing histamine and other mediators within minutes. Examples include hay fever, urticaria, food allergy, and anaphylaxis.
    \item \textbf{Type II (antibody-mediated, cytotoxic):} IgG or IgM antibodies bind to antigens on the surface of the body's own cells, marking them for destruction. Examples include some drug-induced haemolytic anaemia and incompatible blood transfusion reactions.
    \item \textbf{Type III (immune-complex mediated):} circulating antibody-antigen complexes deposit in tissues such as skin, joints, and kidneys, triggering complement activation and inflammation days after exposure. Examples include serum sickness and some drug reactions.
    \item \textbf{Type IV (delayed, cell-mediated):} sensitised T lymphocytes release cytokines and recruit macrophages, producing inflammation hours to days later. Examples include contact dermatitis from nickel or poison ivy, and tuberculin skin reactions.
\end{itemize}

Many real-world reactions, particularly drug allergies, involve more than one mechanism, and the same drug may produce different reaction patterns in different individuals.

\section*{Clinical Features}
Symptoms vary by mechanism and body system:

\begin{itemize}
    \item \textbf{Type I:} sneezing, runny or blocked nose, itchy eyes, urticaria, angio-oedema, wheezing, vomiting, and, in severe cases, anaphylaxis with hypotension and collapse.
    \item \textbf{Type II:} destruction of blood cells, causing anaemia, bruising or bleeding, and low white cell counts, sometimes accompanied by fever and jaundice.
    \item \textbf{Type III:} fever, joint pain, rash, and renal impairment appearing days after exposure to a drug or other antigen.
    \item \textbf{Type IV:} an itching, erythematous, often blistering rash localised to the site of skin contact, developing over 24 to 72 hours.
    \item \textbf{Drug allergy:} can combine features of several mechanisms and may present with widespread rash, fever, organ dysfunction, or the severe forms of Stevens-Johnson syndrome and toxic epidermal necrolysis.
\end{itemize}

\section*{Differential Diagnosis}
Not every adverse reaction is a hypersensitivity reaction:

\begin{itemize}
    \item \textbf{Food intolerance:} including lactose intolerance, with no immune involvement.
    \item \textbf{Drug side effects:} predictable pharmacological effects such as drowsiness from antihistamines, not immune-mediated.
    \item \textbf{Viral infections:} causing rashes, fevers, or joint aches that mimic reactions.
    \item \textbf{Autoimmune diseases:} related in mechanism but distinct in their triggers and chronicity.
    \item \textbf{Non-allergic rhinitis and asthma:} triggered by irritants without an allergic basis.
    \item \textbf{Mast cell disorders:} such as mastocytosis, which cause allergic-type symptoms without allergen exposure.
    \item \textbf{Anxiety and hyperventilation:} which can mimic symptoms of a reaction.
\end{itemize}

\section*{When to Seek Medical Care}
See a doctor if a rash, hives, fever, or other symptoms follow a new medicine, food, or exposure, and if allergic symptoms are frequent, severe, or affecting daily life. Any reaction to a newly prescribed medicine should be reported promptly, because the medicine may need to be stopped or changed.

\textbf{Contact your local emergency services for:}
\begin{itemize}
    \item Difficulty breathing, wheezing, or a hoarse voice.
    \item Swelling of the face, tongue, or throat, or difficulty swallowing.
    \item Dizziness, collapse, or loss of consciousness.
    \item Any signs of anaphylaxis or a severe reaction to a medicine, including widespread blistering rash or skin pain.
\end{itemize}

\section*{Diagnosis and Investigations}
Identifying the cause of a hypersensitivity reaction requires a careful, structured approach:

\begin{itemize}
    \item \textbf{Detailed history:} the substance involved, the timing and nature of the reaction, prior exposures, and any personal or family history of allergy.
    \item \textbf{Clinical examination:} assessment of the skin, mucous membranes, joints, and other affected systems.
    \item \textbf{Skin prick testing and intradermal testing:} for immediate-type allergic reactions.
    \item \textbf{Patch testing:} for delayed contact dermatitis, performed by a dermatologist.
    \item \textbf{Serum specific IgE:} blood tests for allergen-specific antibodies.
    \item \textbf{Haematology:} full blood count and appropriate tests where cell destruction is suspected.
    \item \textbf{Supervised drug or food challenges:} carefully conducted rechallenge in a specialist setting when the diagnosis is uncertain.
    \item \textbf{Referral:} to an allergist, clinical immunologist, or dermatologist for complex or severe cases.
\end{itemize}

\section*{Management}
Treatment depends on the type and severity of the reaction:

\begin{itemize}
    \item \textbf{Strict avoidance:} of the responsible allergen, medicine, or contactant.
    \item \textbf{Antihistamines:} for itching, sneezing, and urticaria.
    \item \textbf{Topical or systemic corticosteroids:} for allergic and immune-mediated inflammatory reactions.
    \item \textbf{Immunosuppressants:} for severe antibody-mediated and immune-complex reactions.
    \item \textbf{Adrenaline:} the first-line treatment for anaphylaxis, delivered by auto-injector.
    \item \textbf{Desensitisation:} gradual reintroduction of an essential medicine under hospital supervision when no alternative exists.
    \item \textbf{Allergen immunotherapy:} for selected inhaled allergens and insect venom.
    \item \textbf{Written action plans and medical alert identification:} for people at risk of severe reactions.
\end{itemize}

\section*{Complications}
Hypersensitivity reactions can have serious consequences:

\begin{itemize}
    \item Anaphylaxis, a medical emergency that can be fatal without prompt treatment.
    \item Severe cutaneous drug reactions, including Stevens-Johnson syndrome and toxic epidermal necrolysis.
    \item Organ damage from immune-complex deposition, affecting the kidneys, joints, or other tissues.
    \item Destruction of blood cells causing anaemia, bleeding, or increased infection risk.
    \item Loss of useful treatments if a drug allergy is wrongly assumed to be present.
    \item Anxiety and restriction of daily activities in people with severe allergies.
\end{itemize}

\section*{Prognosis and Outcomes}
Outcomes vary with the type and severity of the reaction. Most type I allergies are manageable with avoidance and medication, and some, particularly in children, improve with age. Drug reactions usually resolve once the offending medicine is stopped, although a minority are severe or life-threatening. Conditions involving antibody and immune-complex mechanisms may require longer-term specialist care. With accurate diagnosis and an individualised management plan, most people with hypersensitivity disorders lead full and normal lives.

\section*{Prevention and Self-Care}
Practical steps reduce risk and impact:

\begin{itemize}
    \item Record all drug allergies in your medical history and on allergy alerts.
    \item Avoid known allergens and read labels on foods, cosmetics, and medicines.
    \item Carry an adrenaline auto-injector and an up-to-date action plan if prescribed.
    \item Wear medical alert identification if you have had a severe reaction.
    \item Ask about a medicine's ingredients before taking anything new if you have a drug allergy history.
    \item Report any reaction promptly so it can be investigated and documented.
\end{itemize}

\section*{Sources and Further Reading}
The following reputable sources provide reliable, up-to-date information on allergies and hypersensitivity:

\begin{itemize}
    \item NHS. \textit{Allergies and hypersensitivity information for patients.} https://www.nhs.uk/
    \item Mayo Clinic. \textit{Allergies: symptoms, causes, and treatment.} https://www.mayoclinic.org/
    \item MSD Manual. \textit{Overview of allergic and atopic disorders.} https://www.msdmanuals.com/
    \item MedlinePlus. \textit{Allergy and immune system health information.} https://medlineplus.gov/
    \item World Health Organization. \textit{Information on allergic diseases.} https://www.who.int/
    \item NICE. \textit{Clinical guidance on drug allergy and adverse drug reactions.} https://www.nice.org.uk/
\end{itemize}
